\section{common/vecmatop.hpp File Reference}
\label{vecmatop_8hpp}\index{common/vecmatop.hpp@{common/vecmatop.hpp}}


\subsection{Detailed Description}


Definition in file {\bf vecmatop.hpp}.

{\tt \#include \char`\"{}nummat.hpp\char`\"{}}\par
\subsection*{Enumerations}
\begin{CompactItemize}
\item 
enum \{ {\bf EVFLAG\_\-VL} =  1, 
{\bf EVFLAG\_\-FD} =  2, 
{\bf EVFLAG\_\-SD} =  4
 \}
\item 
enum \{ {\bf DMFLAG\_\-PERIOD} =  0, 
{\bf DMFLAG\_\-CLOSED} =  1
 \}
\end{CompactItemize}
\subsection*{Functions}
\begin{CompactItemize}
\item 
int {\bf dscal} (double alpha, {\bf Dbl\-Num\-Vec} \&X)
\item 
int {\bf dscal} (int n, double alpha, double $\ast$X)
\item 
int {\bf daxpy} (double a, const {\bf Dbl\-Num\-Vec} \&X, {\bf Dbl\-Num\-Vec} \&Y)
\item 
int {\bf daxpy} (double a, const {\bf Dbl\-Num\-Mat} \&X, {\bf Dbl\-Num\-Mat} \&Y)
\item 
int {\bf daxpy} (int n, double a, double $\ast$X, double $\ast$Y)
\item 
int {\bf dgemm} (double alpha, const {\bf Dbl\-Num\-Mat} \&A, const {\bf Dbl\-Num\-Mat} \&B, double beta, {\bf Dbl\-Num\-Mat} \&C)
\item 
int {\bf dgemm} (int m, int n, int k, double alpha, double $\ast$A, double $\ast$B, double beta, double $\ast$C)
\item 
int {\bf dger} (double alpha, const {\bf Dbl\-Num\-Vec} \&X, const {\bf Dbl\-Num\-Vec} \&Y, {\bf Dbl\-Num\-Mat} \&A)
\item 
int {\bf dger} (int m, int n, double alpha, double $\ast$X, double $\ast$Y, double $\ast$A)
\item 
int {\bf dgemv} (double alpha, const {\bf Dbl\-Num\-Mat} \&A, const {\bf Dbl\-Num\-Vec} \&X, double beta, {\bf Dbl\-Num\-Vec} \&Y)
\item 
int {\bf dgemv} (int m, int n, double alpha, double $\ast$A, double $\ast$X, double beta, double $\ast$Y)
\item 
int {\bf tran} (const {\bf Dbl\-Num\-Mat} \&M, {\bf Dbl\-Num\-Mat} \&R)
\item 
int {\bf pinv} (const {\bf Dbl\-Num\-Mat} \&M, double epsilon, {\bf Dbl\-Num\-Mat} \&R)
\item 
int {\bf inv} (const {\bf Dbl\-Num\-Mat} \&M, {\bf Dbl\-Num\-Mat} \&R)
\item 
int {\bf spev1d} (int evflag, int dmflag, int dof, double $\ast$M, int n, double e, int i, double u, double $\ast$res)
\item 
int {\bf spev2d} (int evflag, int dmflag, int dof, double $\ast$M, int $\ast$mn, double $\ast$ef, int $\ast$ij, double $\ast$uv, double $\ast$res)
\item 
int {\bf spcoef} (int evflag, double u, double $\ast$us)
\item 
int {\bf lagev1d} (int evflag, int pdflag, int dof, double $\ast$M, int n, double e, int i, double u, double $\ast$res)
\item 
int {\bf lagev2d} (int evflag, int pdflag, int dof, double $\ast$M, int $\ast$mn, double $\ast$ef, int $\ast$ij, double $\ast$uv, double $\ast$res)
\item 
int {\bf lagev3d} (int evflag, int pdflag, int dof, double $\ast$M, int $\ast$mno, double $\ast$efg, int $\ast$ijk, double $\ast$uvw, double $\ast$res)
\item 
int {\bf lagcoef} (int evflag, double u, double $\ast$us)
\end{CompactItemize}


\subsection{Enumeration Type Documentation}
\subsubsection{\setlength{\rightskip}{0pt plus 5cm}anonymous enum}\label{vecmatop_8hpp_dca29a1140aadadfd92b34a02fa516ef_dca29a1140aadadfd92b34a02fa516ef}


\begin{Desc}
\item[Enumerator: ]\par
\begin{description}
\index{EVFLAG_VL@{EVFLAG\_\-VL}!vecmatop.hpp@{vecmatop.hpp}}\index{vecmatop.hpp@{vecmatop.hpp}!EVFLAG_VL@{EVFLAG\_\-VL}}\item[{\em 
EVFLAG\_\-VL\label{vecmatop_8hpp_dca29a1140aadadfd92b34a02fa516ef_dca29a1140aadadfd92b34a02fa516ef8771bdfa1f2367105af9a7eeaf9efb3c}
}]\index{EVFLAG_FD@{EVFLAG\_\-FD}!vecmatop.hpp@{vecmatop.hpp}}\index{vecmatop.hpp@{vecmatop.hpp}!EVFLAG_FD@{EVFLAG\_\-FD}}\item[{\em 
EVFLAG\_\-FD\label{vecmatop_8hpp_dca29a1140aadadfd92b34a02fa516ef_dca29a1140aadadfd92b34a02fa516efacb808ed11f20304d1ee29bc77d807fb}
}]\index{EVFLAG_SD@{EVFLAG\_\-SD}!vecmatop.hpp@{vecmatop.hpp}}\index{vecmatop.hpp@{vecmatop.hpp}!EVFLAG_SD@{EVFLAG\_\-SD}}\item[{\em 
EVFLAG\_\-SD\label{vecmatop_8hpp_dca29a1140aadadfd92b34a02fa516ef_dca29a1140aadadfd92b34a02fa516ef037aa5285a9d7eabc5e9902e5bba2b1d}
}]\end{description}
\end{Desc}



Definition at line 51 of file vecmatop.hpp.

\footnotesize\begin{verbatim}51 {  EVFLAG_VL = 1,  EVFLAG_FD = 2,  EVFLAG_SD = 4 };
\end{verbatim}\normalsize 


\subsubsection{\setlength{\rightskip}{0pt plus 5cm}anonymous enum}\label{vecmatop_8hpp_68986ab776eb5d6b5a809a1c005a7300_68986ab776eb5d6b5a809a1c005a7300}


\begin{Desc}
\item[Enumerator: ]\par
\begin{description}
\index{DMFLAG_PERIOD@{DMFLAG\_\-PERIOD}!vecmatop.hpp@{vecmatop.hpp}}\index{vecmatop.hpp@{vecmatop.hpp}!DMFLAG_PERIOD@{DMFLAG\_\-PERIOD}}\item[{\em 
DMFLAG\_\-PERIOD\label{vecmatop_8hpp_68986ab776eb5d6b5a809a1c005a7300_68986ab776eb5d6b5a809a1c005a7300a8081f7a07a399ed0e6dcb3162291bde}
}]\index{DMFLAG_CLOSED@{DMFLAG\_\-CLOSED}!vecmatop.hpp@{vecmatop.hpp}}\index{vecmatop.hpp@{vecmatop.hpp}!DMFLAG_CLOSED@{DMFLAG\_\-CLOSED}}\item[{\em 
DMFLAG\_\-CLOSED\label{vecmatop_8hpp_68986ab776eb5d6b5a809a1c005a7300_68986ab776eb5d6b5a809a1c005a7300df94bb986d4f0a380c3ae5e7170b71f2}
}]\end{description}
\end{Desc}



Definition at line 53 of file vecmatop.hpp.

\footnotesize\begin{verbatim}53 {  DMFLAG_PERIOD = 0,  DMFLAG_CLOSED = 1 };
\end{verbatim}\normalsize 




\subsection{Function Documentation}
\index{vecmatop.hpp@{vecmatop.hpp}!daxpy@{daxpy}}
\index{daxpy@{daxpy}!vecmatop.hpp@{vecmatop.hpp}}
\subsubsection{\setlength{\rightskip}{0pt plus 5cm}int daxpy (int {\em n}, double {\em a}, double $\ast$ {\em X}, double $\ast$ {\em Y})}\label{vecmatop_8hpp_42f36f4fc8b532e752c7329e335ac044_42f36f4fc8b532e752c7329e335ac044}




Definition at line 57 of file vecmatop.cpp.

References DAXPY.

\footnotesize\begin{verbatim}58 {
59   int incx = 1;  int incy = 1;
60   DAXPY(&n, &a, X, &incx, Y, &incy);
61   return 0;
62 }
\end{verbatim}\normalsize 


\index{vecmatop.hpp@{vecmatop.hpp}!daxpy@{daxpy}}
\index{daxpy@{daxpy}!vecmatop.hpp@{vecmatop.hpp}}
\subsubsection{\setlength{\rightskip}{0pt plus 5cm}int daxpy (double {\em a}, const {\bf Dbl\-Num\-Mat} \& {\em X}, {\bf Dbl\-Num\-Mat} \& {\em Y})}\label{vecmatop_8hpp_9e02e7b9fbd857494e5ee749480735cf_9e02e7b9fbd857494e5ee749480735cf}




Definition at line 50 of file vecmatop.cpp.

References Num\-Mat$<$ F $>$::data(), daxpy(), i\-C, Num\-Mat$<$ F $>$::m(), and Num\-Mat$<$ F $>$::n().

\footnotesize\begin{verbatim}51 {
52   assert( X.m() == Y.m() );  assert( X.n() == Y.n() );
53   iC( daxpy(X.m()*X.n(), a, X.data(), Y.data()) );
54   return 0;
55 }
\end{verbatim}\normalsize 


\index{vecmatop.hpp@{vecmatop.hpp}!daxpy@{daxpy}}
\index{daxpy@{daxpy}!vecmatop.hpp@{vecmatop.hpp}}
\subsubsection{\setlength{\rightskip}{0pt plus 5cm}int daxpy (double {\em a}, const {\bf Dbl\-Num\-Vec} \& {\em X}, {\bf Dbl\-Num\-Vec} \& {\em Y})}\label{vecmatop_8hpp_217995e82dfd0f0d81770835a256a079_217995e82dfd0f0d81770835a256a079}




Definition at line 43 of file vecmatop.cpp.

References Num\-Vec$<$ F $>$::data(), daxpy(), and Num\-Vec$<$ F $>$::m().

\footnotesize\begin{verbatim}44 {
45   assert( X.m() == Y.m() );
46   daxpy(X.m(), a, X.data(), Y.data());
47   return 0;
48 }
\end{verbatim}\normalsize 


\index{vecmatop.hpp@{vecmatop.hpp}!dgemm@{dgemm}}
\index{dgemm@{dgemm}!vecmatop.hpp@{vecmatop.hpp}}
\subsubsection{\setlength{\rightskip}{0pt plus 5cm}int dgemm (int {\em m}, int {\em n}, int {\em k}, double {\em alpha}, double $\ast$ {\em A}, double $\ast$ {\em B}, double {\em beta}, double $\ast$ {\em C})}\label{vecmatop_8hpp_ab0f3ccb6c8a9c6edf137f60264afe2c_ab0f3ccb6c8a9c6edf137f60264afe2c}




Definition at line 72 of file vecmatop.cpp.

References DGEMM.

\footnotesize\begin{verbatim}73 {
74   char transa = 'N';
75   char transb = 'N';
76   assert(m!=0 && n!=0 && k!=0);
77   DGEMM(&transa, &transb, &m, &n, &k,
78                   &alpha, A, &m, B, &k, &beta, C, &m);
79   return 0;
80 }
\end{verbatim}\normalsize 


\index{vecmatop.hpp@{vecmatop.hpp}!dgemm@{dgemm}}
\index{dgemm@{dgemm}!vecmatop.hpp@{vecmatop.hpp}}
\subsubsection{\setlength{\rightskip}{0pt plus 5cm}int dgemm (double {\em alpha}, const {\bf Dbl\-Num\-Mat} \& {\em A}, const {\bf Dbl\-Num\-Mat} \& {\em B}, double {\em beta}, {\bf Dbl\-Num\-Mat} \& {\em C})}\label{vecmatop_8hpp_7893f4c4a19f9a71147f10666feaa39e_7893f4c4a19f9a71147f10666feaa39e}




Definition at line 65 of file vecmatop.cpp.

References Num\-Mat$<$ F $>$::data(), dgemm(), i\-C, Num\-Mat$<$ F $>$::m(), and Num\-Mat$<$ F $>$::n().

\footnotesize\begin{verbatim}66 {
67   assert( A.m() == C.m() );  assert( A.n() == B.m() );  assert( B.n() == C.n() );
68   iC( dgemm(C.m(), C.n(), A.n(), alpha, A.data(), B.data(), beta, C.data()) );
69   return 0;
70 }
\end{verbatim}\normalsize 


\index{vecmatop.hpp@{vecmatop.hpp}!dgemv@{dgemv}}
\index{dgemv@{dgemv}!vecmatop.hpp@{vecmatop.hpp}}
\subsubsection{\setlength{\rightskip}{0pt plus 5cm}int dgemv (int {\em m}, int {\em n}, double {\em alpha}, double $\ast$ {\em A}, double $\ast$ {\em X}, double {\em beta}, double $\ast$ {\em Y})}\label{vecmatop_8hpp_f86104cd2d030678903b8a17738ff742_f86104cd2d030678903b8a17738ff742}




Definition at line 107 of file vecmatop.cpp.

References DGEMV.

\footnotesize\begin{verbatim}108 {
109   char trans = 'N';
110   assert(m!=0 && n!=0);
111   int incx = 1;
112   int incy = 1;
113   DGEMV(&trans, &m, &n, &alpha, A, &m, X, &incx, &beta, Y, &incy);
114   return 0;
115 }
\end{verbatim}\normalsize 


\index{vecmatop.hpp@{vecmatop.hpp}!dgemv@{dgemv}}
\index{dgemv@{dgemv}!vecmatop.hpp@{vecmatop.hpp}}
\subsubsection{\setlength{\rightskip}{0pt plus 5cm}int dgemv (double {\em alpha}, const {\bf Dbl\-Num\-Mat} \& {\em A}, const {\bf Dbl\-Num\-Vec} \& {\em X}, double {\em beta}, {\bf Dbl\-Num\-Vec} \& {\em Y})}\label{vecmatop_8hpp_1f65eba330c9ba99a02f9ec4ecdc51ef_1f65eba330c9ba99a02f9ec4ecdc51ef}




Definition at line 99 of file vecmatop.cpp.

References Num\-Vec$<$ F $>$::data(), Num\-Mat$<$ F $>$::data(), dgemv(), i\-C, Num\-Mat$<$ F $>$::m(), Num\-Vec$<$ F $>$::m(), and Num\-Mat$<$ F $>$::n().

\footnotesize\begin{verbatim}100 {
101   assert(Y.m() == A.m());
102   assert(A.n() == X.m());
103   iC( dgemv(A.m(), A.n(), alpha, A.data(), X.data(), beta, Y.data()) );
104   return 0;
105 }
\end{verbatim}\normalsize 


\index{vecmatop.hpp@{vecmatop.hpp}!dger@{dger}}
\index{dger@{dger}!vecmatop.hpp@{vecmatop.hpp}}
\subsubsection{\setlength{\rightskip}{0pt plus 5cm}int dger (int {\em m}, int {\em n}, double {\em alpha}, double $\ast$ {\em X}, double $\ast$ {\em Y}, double $\ast$ {\em A})}\label{vecmatop_8hpp_f8ede744e4705c3d338ad70b5b4f5eef_f8ede744e4705c3d338ad70b5b4f5eef}




Definition at line 90 of file vecmatop.cpp.

References DGER.

\footnotesize\begin{verbatim}91 {
92   assert(m!=0 && n!=0);
93   int incx = 1;  int incy = 1;
94   DGER(&m, &n, &alpha, X, &incx, Y, &incy, A, &m);
95   return 0;
96 }
\end{verbatim}\normalsize 


\index{vecmatop.hpp@{vecmatop.hpp}!dger@{dger}}
\index{dger@{dger}!vecmatop.hpp@{vecmatop.hpp}}
\subsubsection{\setlength{\rightskip}{0pt plus 5cm}int dger (double {\em alpha}, const {\bf Dbl\-Num\-Vec} \& {\em X}, const {\bf Dbl\-Num\-Vec} \& {\em Y}, {\bf Dbl\-Num\-Mat} \& {\em A})}\label{vecmatop_8hpp_a7e716e5ac5169dc959018df48f6aafa_a7e716e5ac5169dc959018df48f6aafa}




Definition at line 82 of file vecmatop.cpp.

References Num\-Mat$<$ F $>$::data(), Num\-Vec$<$ F $>$::data(), dger(), i\-C, Num\-Mat$<$ F $>$::m(), Num\-Vec$<$ F $>$::m(), and Num\-Mat$<$ F $>$::n().

\footnotesize\begin{verbatim}83 {
84   assert(X.m() == A.m());
85   assert(Y.m() == A.n());
86   iC( dger(A.m(), A.n(), alpha, X.data(), Y.data(), A.data()) );
87   return 0;
88 }
\end{verbatim}\normalsize 


\index{vecmatop.hpp@{vecmatop.hpp}!dscal@{dscal}}
\index{dscal@{dscal}!vecmatop.hpp@{vecmatop.hpp}}
\subsubsection{\setlength{\rightskip}{0pt plus 5cm}int dscal (int {\em n}, double {\em alpha}, double $\ast$ {\em X})}\label{vecmatop_8hpp_7358d64948ad9a47cf0744506e881c1d_7358d64948ad9a47cf0744506e881c1d}




Definition at line 36 of file vecmatop.cpp.

References DSCAL.

\footnotesize\begin{verbatim}37 {
38   int incx = 1;
39   DSCAL(&n, &alpha, X, &incx);
40   return 0;
41 }
\end{verbatim}\normalsize 


\index{vecmatop.hpp@{vecmatop.hpp}!dscal@{dscal}}
\index{dscal@{dscal}!vecmatop.hpp@{vecmatop.hpp}}
\subsubsection{\setlength{\rightskip}{0pt plus 5cm}int dscal (double {\em alpha}, {\bf Dbl\-Num\-Vec} \& {\em X})}\label{vecmatop_8hpp_3b4887d360a4d7483a4a1060d10012f7_3b4887d360a4d7483a4a1060d10012f7}




Definition at line 30 of file vecmatop.cpp.

References Num\-Vec$<$ F $>$::data(), dscal(), and Num\-Vec$<$ F $>$::m().

\footnotesize\begin{verbatim}31 {
32   dscal( X.m(), alpha, X.data() );
33   return 0;
34 }
\end{verbatim}\normalsize 


\index{vecmatop.hpp@{vecmatop.hpp}!inv@{inv}}
\index{inv@{inv}!vecmatop.hpp@{vecmatop.hpp}}
\subsubsection{\setlength{\rightskip}{0pt plus 5cm}int inv (const {\bf Dbl\-Num\-Mat} \& {\em M}, {\bf Dbl\-Num\-Mat} \& {\em R})}\label{vecmatop_8hpp_9b79730e9b5f630d287f00e551e8c344_9b79730e9b5f630d287f00e551e8c344}




Definition at line 164 of file vecmatop.cpp.

References Num\-Mat$<$ F $>$::data(), DGETRF, DGETRI, Num\-Mat$<$ F $>$::m(), and Num\-Mat$<$ F $>$::n().

\footnotesize\begin{verbatim}165 {
166   //OR pinv(M, 0.0, R);
167   assert(M.m()==M.n() && R.m()==R.n() && M.m()==R.m());
168   memcpy(R.data(), M.data(), M.m()*M.n()*sizeof(double));
169   int info;
170   int m = M.m();
171   int* ipiv = new int[m];
172   DGETRF(&m, &m, R.data(), &m, ipiv, &info); assert(info==0);
173   int lwork = m;
174   double* work = new double[lwork];
175   DGETRI(&m, R.data(), &m, ipiv, work, &lwork, &info);  assert(info==0);
176   delete [] ipiv;
177   delete [] work;
178   return 0;
179 }
\end{verbatim}\normalsize 


\index{vecmatop.hpp@{vecmatop.hpp}!lagcoef@{lagcoef}}
\index{lagcoef@{lagcoef}!vecmatop.hpp@{vecmatop.hpp}}
\subsubsection{\setlength{\rightskip}{0pt plus 5cm}int lagcoef (int {\em evflag}, double {\em u}, double $\ast$ {\em us})}\label{vecmatop_8hpp_ae45c5a2ae5bb64731d2e859f5ca0586_ae45c5a2ae5bb64731d2e859f5ca0586}




Definition at line 564 of file vecmatop.cpp.

\footnotesize\begin{verbatim}565 {
566   //double u1 = u;
567   double u2 = u*u;
568   double u3 = u*u*u;
569   if(       evflag==EVFLAG_VL) {
570          us[0] = (u3-3*u2+2*u)/(-6.0);
571          us[1] = (u3-2*u2-u+2)/(2.0);
572          us[2] = (u3-u2-2*u)  /(-2.0);
573          us[3] = (u3-u)       /(6.0);
574   } else if(evflag==EVFLAG_FD) {
575          us[0] = (3*u2-6*u+2)/(-6.0);
576          us[1] = (3*u2-4*u-1)/(2.0);
577          us[2] = (3*u2-2*u-2)/(-2.0);
578          us[3] = (3*u2-1)    /(6.0);
579   } else if(evflag==EVFLAG_SD) {
580          assert(0); //TODO
581   }
582   return 0;
583 }
\end{verbatim}\normalsize 


\index{vecmatop.hpp@{vecmatop.hpp}!lagev1d@{lagev1d}}
\index{lagev1d@{lagev1d}!vecmatop.hpp@{vecmatop.hpp}}
\subsubsection{\setlength{\rightskip}{0pt plus 5cm}int lagev1d (int {\em evflag}, int {\em pdflag}, int {\em dof}, double $\ast$ {\em M}, int {\em n}, double {\em e}, int {\em i}, double {\em u}, double $\ast$ {\em res})}\label{vecmatop_8hpp_559daf9c86280b25416fa70e8a206c82_559daf9c86280b25416fa70e8a206c82}




Definition at line 363 of file vecmatop.cpp.

References DMFLAG\_\-PERIOD.

\footnotesize\begin{verbatim}364 {
365   int is[4];
366   if(dmflag==DMFLAG_PERIOD) {
367          for(int k=0; k<4; k++) is[k]=(i+k-1 + m) % m;
368   } else {
369          assert(i>=1 && i<=m-3);
370          for(int k=0; k<4; k++) is[k]=(i+k-1);
371   }
372   DblNumMat M(dof,m,false,data); //assert(M.n()==n && M.m()==res.m()); //double dof = M.m();  //int cnt = 0;
373   //---------------------------
374   if(evflag & EVFLAG_VL) {
375          double scl = 1.0;
376          double us[4]; iC( lagcoef(EVFLAG_VL, u, us) );
377          for(int d=0; d<dof; d++)               res[d] = 0;
378          for(int d=0; d<dof; d++) 
379                 for(int a=0; a<4; a++) {
380                   res[d] += us[a] * M(d,is[a]);
381                 }
382          for(int d=0; d<dof; d++)               res[d] *= scl; //scaling
383          res+=dof; //cnt ++;
384   }
385   //---------------------------
386   if(evflag & EVFLAG_FD) {
387          double scl = double(m) / e;
388          double us[4]; iC( lagcoef(EVFLAG_FD, u, us) );
389          for(int d=0; d<dof; d++)               res[d] = 0;
390          for(int d=0; d<dof; d++)
391                 for(int a=0; a<4; a++) {
392                   res[d] += us[a] * M(d,is[a]);
393                 }
394          for(int d=0; d<dof; d++)               res[d] *= scl; //scaling
395          res+=dof; //cnt ++;
396   }
397   //---------------------------
398   if(evflag & EVFLAG_SD) {
399          double scl = double(m*m)/(e*e);
400          double us[4]; iC( lagcoef(EVFLAG_SD, u, us) );
401          for(int d=0; d<dof; d++)               res[d] = 0;
402          for(int d=0; d<dof; d++)
403                 for(int a=0; a<4; a++) {
404                   res[d] += us[a] * M(d,is[a]);
405                 }
406          for(int d=0; d<dof; d++)               res[d] *= scl; //scaling
407          res+=dof; //cnt ++;
408   }
409   return 0;
410 }
\end{verbatim}\normalsize 


\index{vecmatop.hpp@{vecmatop.hpp}!lagev2d@{lagev2d}}
\index{lagev2d@{lagev2d}!vecmatop.hpp@{vecmatop.hpp}}
\subsubsection{\setlength{\rightskip}{0pt plus 5cm}int lagev2d (int {\em evflag}, int {\em pdflag}, int {\em dof}, double $\ast$ {\em M}, int $\ast$ {\em mn}, double $\ast$ {\em ef}, int $\ast$ {\em ij}, double $\ast$ {\em uv}, double $\ast$ {\em res})}\label{vecmatop_8hpp_5777432f838d73947c2307cc63592160_5777432f838d73947c2307cc63592160}




Definition at line 412 of file vecmatop.cpp.

References DMFLAG\_\-PERIOD.

\footnotesize\begin{verbatim}413 {
414   int m = mn[0];  int n = mn[1];
415   double e = ef[0];  double f = ef[1];
416   int i = ij[0];  int j = ij[1];
417   double u = uv[0];  double v = uv[1];
418   
419   int is[4]; int js[4];
420   if(dmflag==DMFLAG_PERIOD) {
421          for(int k=0; k<4; k++) is[k]=(i+k-1 + m) % m;
422          for(int k=0; k<4; k++) js[k]=(j+k-1 + n) % n;
423   } else {
424          assert(i>=1 && i<=m-3);
425          for(int k=0; k<4; k++) is[k]=(i+k-1);
426          assert(j>=1 && j<=n-3);
427          for(int k=0; k<4; k++) js[k]=(j+k-1);
428   }
429   DblNumMat M(dof,m*n,false,data);
430   double scl; 
431   double us[4], vs[4];
432   //---------------------------
433   if(evflag & EVFLAG_VL) {
434          scl = 1.0;
435          iC( lagcoef(EVFLAG_VL, u, us) );
436          iC( lagcoef(EVFLAG_VL, v, vs) );
437          for(int d=0; d<dof; d++)               res[d] = 0;
438          for(int a=0; a<4; a++)
439                 for(int b=0; b<4; b++) {
440                   double coef = us[a]*vs[b]; 
441                   for(int d=0; d<dof; d++)
442                          res[d] += coef * M(d, is[a]+js[b]*m);
443                 }
444          for(int d=0; d<dof; d++)               res[d] *= scl;
445          res+=dof;
446   }
447   //---------------------------
448   if(evflag & EVFLAG_FD) {
449          scl = double(m)/e;
450          iC( lagcoef(EVFLAG_FD, u, us) );
451          iC( lagcoef(EVFLAG_VL, v, vs) );
452          for(int d=0; d<dof; d++)               res[d] = 0;
453          for(int a=0; a<4; a++)
454                 for(int b=0; b<4; b++) {
455                   double coef = us[a]*vs[b];
456                   for(int d=0; d<dof; d++)
457                          res[d] += coef * M(d, is[a]+js[b]*m);
458                 }
459          for(int d=0; d<dof; d++)               res[d] *= scl;
460          res+=dof;
461          //...
462          scl = double(n)/f;
463          iC( lagcoef(EVFLAG_VL, u, us) );
464          iC( lagcoef(EVFLAG_FD, v, vs) );
465          for(int d=0; d<dof; d++)               res[d] = 0;
466          for(int a=0; a<4; a++)
467                 for(int b=0; b<4; b++) {
468                   double coef = us[a]*vs[b]; 
469                   for(int d=0; d<dof; d++)
470                          res[d] += coef * M(d, is[a]+js[b]*m);
471                 }
472          for(int d=0; d<dof; d++)               res[d] *= scl;
473          res+=dof;
474   }
475   //---------------------------
476   if(evflag & EVFLAG_SD) {
477          scl = double(m*m)/(e*e);
478          iC( lagcoef(EVFLAG_SD, u, us) );
479          iC( lagcoef(EVFLAG_VL, v, vs) );
480          for(int d=0; d<dof; d++)               res[d] = 0;
481          for(int a=0; a<4; a++)
482                 for(int b=0; b<4; b++) {
483                   double coef = us[a]*vs[b]; 
484                   for(int d=0; d<dof; d++)
485                          res[d] += coef * M(d, is[a]+js[b]*m);
486                 }
487          for(int d=0; d<dof; d++)               res[d] *= scl;
488          res+=dof;
489          //...
490          scl = double(m*n)/(e*f);
491          iC( lagcoef(EVFLAG_FD, u, us) );
492          iC( lagcoef(EVFLAG_FD, v, vs) );
493          for(int d=0; d<dof; d++)               res[d] = 0;
494          for(int a=0; a<4; a++)
495                 for(int b=0; b<4; b++) {
496                   double coef = us[a]*vs[b]; 
497                   for(int d=0; d<dof; d++)
498                          res[d] += coef * M(d, is[a]+js[b]*m);
499                 }
500          for(int d=0; d<dof; d++)               res[d] *= scl;
501          res+=dof;
502          //...
503          scl = double(n*n)/(f*f);
504          iC( lagcoef(EVFLAG_VL, u, us) );
505          iC( lagcoef(EVFLAG_SD, v, vs) );
506          for(int d=0; d<dof; d++)               res[d] = 0;
507          for(int a=0; a<4; a++)
508                 for(int b=0; b<4; b++) {
509                   double coef = us[a]*vs[b]; 
510                   for(int d=0; d<dof; d++)
511                          res[d] += coef * M(d, is[a]+js[b]*m);
512                 }
513          for(int d=0; d<dof; d++)               res[d] *= scl;
514          res+=dof;
515   }
516   return 0;
517 }
\end{verbatim}\normalsize 


\index{vecmatop.hpp@{vecmatop.hpp}!lagev3d@{lagev3d}}
\index{lagev3d@{lagev3d}!vecmatop.hpp@{vecmatop.hpp}}
\subsubsection{\setlength{\rightskip}{0pt plus 5cm}int lagev3d (int {\em evflag}, int {\em pdflag}, int {\em dof}, double $\ast$ {\em M}, int $\ast$ {\em mno}, double $\ast$ {\em efg}, int $\ast$ {\em ijk}, double $\ast$ {\em uvw}, double $\ast$ {\em res})}\label{vecmatop_8hpp_505b660797481226ce76fbe2b50d966f_505b660797481226ce76fbe2b50d966f}




Definition at line 519 of file vecmatop.cpp.

References DMFLAG\_\-PERIOD.

\footnotesize\begin{verbatim}520 {
521   assert(evflag==EVFLAG_VL);
522   
523   int m = mno[0];  int n = mno[1];  int o = mno[2];
524   double e = efg[0];  double f = efg[1];  double g = efg[2];
525   int i = ijk[0];  int j = ijk[1];  int k = ijk[2];
526   double u = uvw[0];  double v = uvw[1];  double w = uvw[2];
527   
528   int is[4]; int js[4];  int ks[4];
529   if(dmflag==DMFLAG_PERIOD) {
530          for(int h=0; h<4; h++) is[h]=(i+h-1 + m) % m;
531          for(int h=0; h<4; h++) js[h]=(j+h-1 + n) % n;
532          for(int h=0; h<4; h++) ks[h]=(k+h-1 + o) % o;
533   } else {
534          assert(i>=1 && i<=m-3);
535          for(int h=0; h<4; h++) is[h]=(i+h-1);
536          assert(j>=1 && j<=n-3);
537          for(int h=0; h<4; h++) js[h]=(j+h-1);
538          assert(k>=1 && k<=o-3);
539          for(int h=0; h<4; h++) ks[h]=(k+h-1);
540   }
541   DblNumMat M(dof,m*n*o,false,data);
542   double scl; 
543   double us[4], vs[4], ws[4];
544   //---------------------------
545   if(evflag & EVFLAG_VL) {
546          scl = 1.0;
547          iC( lagcoef(EVFLAG_VL, u, us) );
548          iC( lagcoef(EVFLAG_VL, v, vs) );
549          iC( lagcoef(EVFLAG_VL, w, ws) );
550          for(int d=0; d<dof; d++)               res[d] = 0;
551          for(int a=0; a<4; a++)
552                 for(int b=0; b<4; b++)
553                   for(int c=0; c<4; c++) {
554                          double coef = us[a]*vs[b]*ws[c]; 
555                          for(int d=0; d<dof; d++)
556                                 res[d] += coef * M(d, is[a]+js[b]*m+ks[c]*m*n);
557                   }
558          for(int d=0; d<dof; d++)               res[d] *= scl;
559          res+=dof;
560   }
561   return 0;
562 }
\end{verbatim}\normalsize 


\index{vecmatop.hpp@{vecmatop.hpp}!pinv@{pinv}}
\index{pinv@{pinv}!vecmatop.hpp@{vecmatop.hpp}}
\subsubsection{\setlength{\rightskip}{0pt plus 5cm}int pinv (const {\bf Dbl\-Num\-Mat} \& {\em M}, double {\em epsilon}, {\bf Dbl\-Num\-Mat} \& {\em R})}\label{vecmatop_8hpp_e6044a83323bd043a807de27d79c4365_e6044a83323bd043a807de27d79c4365}




Definition at line 126 of file vecmatop.cpp.

References i\-C, Num\-Mat$<$ F $>$::m(), and Num\-Mat$<$ F $>$::n().

\footnotesize\begin{verbatim}127 {
128   assert(M.m() == R.n());  assert(M.n() == R.m());
129   SVDRep svd;
130   iC( svd.construct(epsilon, M) );
131   //invert Svd    //cerr<<svd.S()(0)<<" "<<svd.S()(svd.S().m()-1)<<" "<<svd.S()(0)/svd.S()(svd.S().m()-1)<<endl;
132   double cutoff = epsilon * svd.S()(0);  //double cutoff = epsilon;
133   for(int i=0; i<svd.S().m(); i++) {
134     if( svd.S()(i) >= cutoff) {
135       svd.S()(i) = 1.0/(svd.S()(i));
136          } else {
137                 assert(0);      //svd.S()(i) = 0.0;
138          }
139   }
140   DblNumMat UT(svd.U().n(),  svd.U().m());
141   DblNumMat V( svd.VT().n(), svd.VT().m());
142   iC( tran(svd.U(), UT) );
143   iC( tran(svd.VT(), V) );
144   for(int i=0; i<V.m(); i++)
145     for(int j=0; j<V.n(); j++) {
146       V(i,j) = V(i,j) * svd.S()(j);
147          }
148   
149   char transa = 'N';
150   char transb = 'N';
151   double alpha = 1.0;
152   double beta = 0.0;
153   int m = V.m();
154   int n = UT.n();
155   int k = V.n();
156   DGEMM(&transa, &transb, &m, &n, &k, &alpha,
157                   V.data(), &m, UT.data(), &k, 
158                   &beta, R.data(), &m);  
159   
160   return 0;
161 }
\end{verbatim}\normalsize 


\index{vecmatop.hpp@{vecmatop.hpp}!spcoef@{spcoef}}
\index{spcoef@{spcoef}!vecmatop.hpp@{vecmatop.hpp}}
\subsubsection{\setlength{\rightskip}{0pt plus 5cm}int spcoef (int {\em evflag}, double {\em u}, double $\ast$ {\em us})}\label{vecmatop_8hpp_9084d4f8c569a1ae492c9b1f315df7c1_9084d4f8c569a1ae492c9b1f315df7c1}




Definition at line 338 of file vecmatop.cpp.

\footnotesize\begin{verbatim}339 {
340   double u1 = u;
341   double u2 = u*u;
342   double u3 = u*u*u;
343   if(       evflag==EVFLAG_VL) {
344          us[0] = (  1 - 3*u1 + 3*u2 -   u3)/6.0;
345          us[1] = (  4        - 6*u2 + 3*u3)/6.0;
346          us[2] = (  1 + 3*u1 + 3*u2 - 3*u3)/6.0;
347          us[3] = (                  +   u3)/6.0;
348   } else if(evflag==EVFLAG_FD) {
349          us[0] = (- 3 + 6*u1 - 3*u2)/6.0;
350          us[1] = (    -12*u1 + 9*u2)/6.0;
351          us[2] = (  3 + 6*u1 - 9*u2)/6.0;
352          us[3] = (             3*u2)/6.0;
353   } else if(evflag==EVFLAG_SD) {
354          us[0] = (  6 - 6*u1 ) / 6.0;
355          us[1] = (-12 +18*u1 ) / 6.0;
356          us[2] = (  6 -18*u1 ) / 6.0;
357          us[3] = (      6*u1 ) / 6.0;    //assert(0); //TODO;
358   }  
359   return 0;
360 }
\end{verbatim}\normalsize 


\index{vecmatop.hpp@{vecmatop.hpp}!spev1d@{spev1d}}
\index{spev1d@{spev1d}!vecmatop.hpp@{vecmatop.hpp}}
\subsubsection{\setlength{\rightskip}{0pt plus 5cm}int spev1d (int {\em evflag}, int {\em dmflag}, int {\em dof}, double $\ast$ {\em M}, int {\em n}, double {\em e}, int {\em i}, double {\em u}, double $\ast$ {\em res})}\label{vecmatop_8hpp_59ca1776ee301a72ccfb5105b3a3edfd_59ca1776ee301a72ccfb5105b3a3edfd}




Definition at line 182 of file vecmatop.cpp.

References DMFLAG\_\-PERIOD.

\footnotesize\begin{verbatim}183 {
184   int is[4];
185   if(dmflag==DMFLAG_PERIOD) {
186          for(int k=0; k<4; k++) is[k]=(i+k-1 + m) % m;
187   } else {
188          assert(i>=1 && i<=m-3);
189          for(int k=0; k<4; k++) is[k]=(i+k-1);
190   }
191   DblNumMat M(dof,m,false,data); //assert(M.n()==n && M.m()==res.m()); //double dof = M.m();  //int cnt = 0;
192   //---------------------------
193   if(evflag & EVFLAG_VL) {
194          double scl = 1.0;
195          double us[4]; iC( spcoef(EVFLAG_VL, u, us) );
196          for(int d=0; d<dof; d++)               res[d] = 0;
197          for(int d=0; d<dof; d++) 
198                 for(int a=0; a<4; a++) {
199                   res[d] += us[a] * M(d,is[a]);
200                 }
201          for(int d=0; d<dof; d++)               res[d] *= scl; //scaling
202          res+=dof;
203   }
204   //---------------------------
205   if(evflag & EVFLAG_FD) {
206          double scl = double(m) / e;
207          double us[4]; iC( spcoef(EVFLAG_FD, u, us) );
208          for(int d=0; d<dof; d++)               res[d] = 0;
209          for(int d=0; d<dof; d++)
210                 for(int a=0; a<4; a++) {
211                   res[d] += us[a] * M(d,is[a]);
212                 }
213          for(int d=0; d<dof; d++)               res[d] *= scl; //scaling
214          res+=dof;
215   }
216   //---------------------------
217   if(evflag & EVFLAG_SD) {
218          double scl = double(m*m)/(e*e);
219          double us[4]; iC( spcoef(EVFLAG_SD, u, us) );
220          for(int d=0; d<dof; d++)               res[d] = 0;
221          for(int d=0; d<dof; d++)
222                 for(int a=0; a<4; a++) {
223                   res[d] += us[a] * M(d,is[a]);
224                 }
225          for(int d=0; d<dof; d++)               res[d] *= scl; //scaling
226          res+=dof;
227   }
228   return 0;
229 }
\end{verbatim}\normalsize 


\index{vecmatop.hpp@{vecmatop.hpp}!spev2d@{spev2d}}
\index{spev2d@{spev2d}!vecmatop.hpp@{vecmatop.hpp}}
\subsubsection{\setlength{\rightskip}{0pt plus 5cm}int spev2d (int {\em evflag}, int {\em dmflag}, int {\em dof}, double $\ast$ {\em M}, int $\ast$ {\em mn}, double $\ast$ {\em ef}, int $\ast$ {\em ij}, double $\ast$ {\em uv}, double $\ast$ {\em res})}\label{vecmatop_8hpp_2bd003f61d68382b102b1f7747d42016_2bd003f61d68382b102b1f7747d42016}




Definition at line 231 of file vecmatop.cpp.

References DMFLAG\_\-PERIOD.

\footnotesize\begin{verbatim}232 {
233   int m = mn[0];  int n = mn[1];
234   double e = ef[0];  double f = ef[1];
235   int i = ij[0];  int j = ij[1];
236   double u = uv[0];  double v = uv[1];
237   
238   int is[4]; int js[4];
239   if(dmflag==DMFLAG_PERIOD) {
240          for(int k=0; k<4; k++) is[k]=(i+k-1 + m) % m;
241          for(int k=0; k<4; k++) js[k]=(j+k-1 + n) % n;
242   } else {
243          assert(i>=1 && i<=m-3);
244          for(int k=0; k<4; k++) is[k]=(i+k-1);
245          assert(j>=1 && j<=n-3);
246          for(int k=0; k<4; k++) js[k]=(j+k-1);
247   }
248   DblNumMat M(dof,m*n,false,data);
249   double scl;
250   double us[4], vs[4];
251   //---------------------------
252   if(evflag & EVFLAG_VL) {
253          scl = 1.0;
254          iC( spcoef(EVFLAG_VL, u, us) );
255          iC( spcoef(EVFLAG_VL, v, vs) );
256          for(int d=0; d<dof; d++)               res[d] = 0;
257          for(int a=0; a<4; a++)
258                 for(int b=0; b<4; b++) {
259                   double coef = us[a]*vs[b]; 
260                   for(int d=0; d<dof; d++)
261                          res[d] += coef * M(d, is[a]+js[b]*m);
262                 }
263          for(int d=0; d<dof; d++)               res[d] *= scl;
264          res+=dof;
265   }
266   //---------------------------
267   if(evflag & EVFLAG_FD) {
268          scl = double(m)/e;
269          iC( spcoef(EVFLAG_FD, u, us) );
270          iC( spcoef(EVFLAG_VL, v, vs) );
271          for(int d=0; d<dof; d++)               res[d] = 0;
272          for(int a=0; a<4; a++)
273                 for(int b=0; b<4; b++) {
274                   double coef = us[a]*vs[b];
275                   for(int d=0; d<dof; d++)
276                          res[d] += coef * M(d, is[a]+js[b]*m);
277                 }
278          for(int d=0; d<dof; d++)               res[d] *= scl;
279          res+=dof;
280          //...
281          scl = double(n)/f;
282          iC( spcoef(EVFLAG_VL, u, us) );
283          iC( spcoef(EVFLAG_FD, v, vs) );
284          for(int d=0; d<dof; d++)               res[d] = 0;
285          for(int a=0; a<4; a++)
286                 for(int b=0; b<4; b++) {
287                   double coef = us[a]*vs[b]; 
288                   for(int d=0; d<dof; d++)
289                          res[d] += coef * M(d, is[a]+js[b]*m);
290                 }
291          for(int d=0; d<dof; d++)               res[d] *= scl;
292          res+=dof;
293   }
294   //---------------------------
295   if(evflag & EVFLAG_SD) {
296          scl = double(m*m)/(e*e);
297          iC( spcoef(EVFLAG_SD, u, us) );
298          iC( spcoef(EVFLAG_VL, v, vs) );
299          for(int d=0; d<dof; d++)               res[d] = 0;
300          for(int a=0; a<4; a++)
301                 for(int b=0; b<4; b++) {
302                   double coef = us[a]*vs[b]; 
303                   for(int d=0; d<dof; d++)
304                          res[d] += coef * M(d, is[a]+js[b]*m);
305                 }
306          for(int d=0; d<dof; d++)               res[d] *= scl;
307          res+=dof;
308          //...
309          scl = double(m*n)/(e*f);
310          iC( spcoef(EVFLAG_FD, u, us) );
311          iC( spcoef(EVFLAG_FD, v, vs) );
312          for(int d=0; d<dof; d++)               res[d] = 0;
313          for(int a=0; a<4; a++)
314                 for(int b=0; b<4; b++) {
315                   double coef = us[a]*vs[b]; 
316                   for(int d=0; d<dof; d++)
317                          res[d] += coef * M(d, is[a]+js[b]*m);
318                 }
319          for(int d=0; d<dof; d++)               res[d] *= scl;
320          res+=dof;
321          //...
322          scl = double(n*n)/(f*f);
323          iC( spcoef(EVFLAG_VL, u, us) );
324          iC( spcoef(EVFLAG_SD, v, vs) );
325          for(int d=0; d<dof; d++)               res[d] = 0;
326          for(int a=0; a<4; a++)
327                 for(int b=0; b<4; b++) {
328                   double coef = us[a]*vs[b]; 
329                   for(int d=0; d<dof; d++)
330                          res[d] += coef * M(d, is[a]+js[b]*m);
331                 }
332          for(int d=0; d<dof; d++)               res[d] *= scl;
333          res+=dof;
334   }
335   return 0;
336 }
\end{verbatim}\normalsize 


\index{vecmatop.hpp@{vecmatop.hpp}!tran@{tran}}
\index{tran@{tran}!vecmatop.hpp@{vecmatop.hpp}}
\subsubsection{\setlength{\rightskip}{0pt plus 5cm}int tran (const {\bf Dbl\-Num\-Mat} \& {\em M}, {\bf Dbl\-Num\-Mat} \& {\em R})}\label{vecmatop_8hpp_2db1bb0aad5cb3c8b5155fb7f23d7b38_2db1bb0aad5cb3c8b5155fb7f23d7b38}




Definition at line 117 of file vecmatop.cpp.

References Num\-Mat$<$ F $>$::m(), and Num\-Mat$<$ F $>$::n().

\footnotesize\begin{verbatim}118 {
119   assert(R.m()==M.n() && R.n()==M.m());  //R.resize(M.n(), M.m());
120   for(int i=0; i<M.m(); i++)
121          for(int j=0; j<M.n(); j++)
122                 R(j,i) = M(i,j);
123   return 0;
124 }
\end{verbatim}\normalsize 


